\documentclass[11pt,a4paper]{article}
\usepackage[utf8]{inputenc}
\usepackage[T1]{fontenc}
\usepackage[french]{babel}
\usepackage{amsmath, amssymb}
\usepackage{siunitx}
\usepackage{tcolorbox}
\usepackage{geometry}
\usepackage{tikz}
\usepackage{xcolor}
\usepackage{fancyhdr}

\geometry{margin=2cm}
\pagestyle{fancy}
\fancyhf{}
\lhead{\textbf{Cahier de Vacances --- R\'{e}visions}}
\rhead{\thepage}

\definecolor{energyblue}{RGB}{30,80,160}
\definecolor{powerred}{RGB}{180,50,50}

\tcbset{sharp corners, colback=white, colframe=black!60, boxrule=0.8pt, left=6pt, right=6pt, top=4pt, bottom=4pt}

\newcommand{\rappelbox}[2]{%
\begin{tcolorbox}[colback=#1!10!white,colframe=#1!60!black,title=\textbf{Rappel}]
#2
\end{tcolorbox}
}

\newcommand{\gridspace}{%
\vspace{0.5cm}
\begin{tikzpicture}[scale=0.5]
\draw[step=0.5cm,gray!30,very thin] (0,0) grid (27,8);
\end{tikzpicture}
\vspace{0.5cm}
}

\begin{document}

\begin{center}
{\LARGE \textbf{Cahier de Vacances}}\\[0.3cm]
{\large R\'{e}visions --- Math\'{e}matiques et Physique}\\[0.2cm]
{\normalsize Reprends les bases \`{a} ton rythme !}
\end{center}

\vspace{1cm}

%% ============================================================
%% CHAPITRE 1
%% ============================================================
\section*{Chapitre 1 -- Nombres relatifs et priorit\'{e}s}

\rappelbox{blue}{%
\begin{itemize}
  \item \textbf{M\^{e}me signe :} on additionne les valeurs absolues et on garde le signe.\\
        Exemple : $(-3)+(-5)=-8$
  \item \textbf{Signes diff\'{e}rents :} on soustrait les valeurs absolues et on garde le signe de la plus grande.\\
        Exemple : $(+8)+(-3)=+5$
  \item \textbf{Soustraction :} soustraire revient \`{a} ajouter l'oppos\'{e} : $a - b = a + (-b)$
  \item \textbf{Priorit\'{e}s :} Parenth\`{e}ses $\rightarrow$ Puissances $\rightarrow$ $\times / \div$ $\rightarrow$ $+ / -$
\end{itemize}
}

\vspace{0.5cm}

\textbf{Exercice 1.} \quad $(-11) + (+7) = $

\gridspace

\textbf{Exercice 2.} \quad $(3) + (-7) = $

\gridspace

\textbf{Exercice 4.} \quad $(-2) - (-9) = $

\gridspace

\newpage

%% ============================================================
%% CHAPITRE 2
%% ============================================================
\section*{Chapitre 2 -- Multiplication et distributivit\'{e}}

\rappelbox{blue}{%
\begin{itemize}
  \item \textbf{M\^{e}mes signes} $\rightarrow$ r\'{e}sultat positif ; \textbf{signes diff\'{e}rents} $\rightarrow$ r\'{e}sultat n\'{e}gatif
  \item \textbf{Distributivit\'{e} :} $a(b+c) = ab + ac$
  \item \textbf{Identit\'{e} remarquable :} $(a+b)^{2} = a^{2} + 2ab + b^{2}$
  \item \textbf{Puissance :} $(a)^{4} = a*a*a*a$
\end{itemize}
}

\vspace{0.5cm}

\textbf{Exercice 1.} \quad $(-3) \times (+4) = $

\gridspace

\textbf{Exercice 3.} \quad D\'{e}velopper : $4 \times (3 + 5) = $

\gridspace

 

\textbf{Exercice 4.} \quad D\'{e}velopper : $(x + 3)^2 = $

\gridspace

\textbf{Exercice 5.} \quad D\'{e}velopper avec l'identit\'{e} remarquable : $(x + y)^{2} = $

\gridspace


\textbf{Exercice 6.} \quad $(3 + 4)(3 - 4) = $

\gridspace

\textbf{Exercice 7.} \quad $(x + y)(x - y) = $

\gridspace

\textbf{Exercice 8.} \quad $\frac{(x + y)(x - y)(x - y)}{(x + y)} = $

\gridspace

\newpage

%% ============================================================
%% CHAPITRE 3
%% ============================================================
\section*{Chapitre 3 -- Fractions}

\rappelbox{green}{%
\begin{itemize}
  \item \textbf{M\^{e}me d\'{e}nominateur :} $\dfrac{a}{b} + \dfrac{c}{b} = \dfrac{a+c}{b}$
  \item \textbf{D\'{e}nominateurs diff\'{e}rents :} trouver un d\'{e}nominateur commun
  \item \textbf{Multiplication :} $\dfrac{a}{b} \times \dfrac{c}{d} = \dfrac{ac}{bd}$
  \item \textbf{Division :} $\dfrac{a}{b} \div \dfrac{c}{d} = \dfrac{a}{b} \times \dfrac{d}{c}$
  \item \textbf{Simplification :} $\dfrac{ac}{bc} = \dfrac{a}{b}$
\end{itemize}
}

\vspace{0.5cm}

\textbf{Exercice 1.} \quad $\dfrac{3}{7} + \dfrac{2}{7} = $

\gridspace

\textbf{Exercice 2.} \quad $\dfrac{1}{3} + \dfrac{1}{4} = $

\gridspace

\textbf{Exercice 3.} \quad $\dfrac{2}{5} \times \dfrac{3}{4} = $

\gridspace

\newpage 

\textbf{Exercice 4.} \quad $\dfrac{3}{4} \div \dfrac{2}{3} = $

\gridspace

\textbf{Exercice 5.} \quad Simplifier $\dfrac{12}{18}$

\gridspace

\textbf{Exercice 6.} \quad $\dfrac{1}{2} + \dfrac{2}{3} - \dfrac{1}{6} = $

\gridspace

\textbf{Exercice 7.} \quad Simplifier en une seule fraction $\dfrac{xy}{x} + \dfrac{xz}{z} - \dfrac{yz}{y} = $

\gridspace

\newpage

%% ============================================================
%% CHAPITRE 4
%% ============================================================
\section*{Chapitre 4 -- Puissances}

\rappelbox{orange}{%
\begin{itemize}
  \item $a^{m} \times a^{n} = a^{m+n}$
  \item $\dfrac{a^{m}}{a^{n}} = a^{m-n}$
  \item $(a^{m})^{n} = a^{m \times n}$
  \item $(ab)^{n} = a^{n}\,b^{n}$
  \item $a^{-n} = \dfrac{1}{a^{n}}$
  \item $a^{0} = 1$
  \item $a^{1/n} = \sqrt[n]{a}$
\end{itemize}
}

\vspace{0.5cm}

\textbf{Exercice 1.} \quad $3^{2} \times 3^{3} = $

\gridspace

\textbf{Exercice 2.} \quad $\dfrac{5^{6}}{5^{4}} = $

\gridspace

\textbf{Exercice 3.} \quad $(2^{3})^{2} = $

\gridspace

\newpage 

\textbf{Exercice 4.} \quad $(3 \times 4)^{2} = $

\gridspace

\textbf{Exercice 5.} \quad $10^{-2} = $

\gridspace

\textbf{Exercice 6.} \quad $7^{0} = $

\gridspace

\textbf{Exercice 7.} \quad $\left(\dfrac{xy}{3z}\right)^{3} \times \left(\dfrac{1}{zyx}\right)^{2} = $

\gridspace


\newpage

%% ============================================================
%% CHAPITRE 5
%% ============================================================
\section*{Chapitre 5 -- Unit\'{e}s SI et conversions}

\rappelbox{blue}{%
\begin{itemize}
  \item \textbf{7 unit\'{e}s de base :} m, kg, s, A, K, mol, cd
  \item \textbf{Pr\'{e}fixes :} k $= 10^{3}$, \; c $= 10^{-2}$, \; m $= 10^{-3}$, \; M $= 10^{6}$, \; $\mu = 10^{-6}$
  \item \textbf{Surfaces :} 2 z\'{e}ros par \'{e}chelon \quad / \quad \textbf{Volumes :} 3 z\'{e}ros par \'{e}chelon
  \item convertir en unité SI avant d'utiliser une formule, il faut assurer la cohérence des unités
  \item $1\,\text{L} = 10^{-3}\,\text{m}^{3}$ \qquad $1\,\text{km/h} = \dfrac{1000}{3600}\,\text{m/s}$ \qquad $1\,\text{bar} = 10^{5}\,\text{Pa}$
\end{itemize}
}

\vspace{0.5cm}

\textbf{Exercice 1.} \quad Convertir $3{,}5\,\text{km}$ en m.

\gridspace

\textbf{Exercice 2.} \quad Convertir $250\,\text{g}$ en kg.

\gridspace

\textbf{Exercice 3.} \quad Convertir $2\,\text{h}$ en s.

\gridspace

\newpage 


\textbf{Exercice 4.} \quad Convertir $36\,\text{km/h}$ en m/s.

\gridspace

\textbf{Exercice 5.} \quad Convertir $50\,\text{cm}^{2}$ en m$^{2}$.

\gridspace

\textbf{Exercice 6.} \quad Simplifier $\dfrac{\text{kg} \cdot \text{m}}{\text{s}^{2}}$ (quelle unit\'{e} ?)

\gridspace

\textbf{Exercice 7.} \quad Simplifier $\dfrac{\text{J}}{\text{s}}$ (quelle unit\'{e} ?)

\gridspace

\newpage

\textbf{Exercice 8.} \quad \text{Convertir 35 [$m^3$] d'eau en [kg]}

\gridspace

\textbf{Exercice 9.} \quad \text{Convertir 27 [$m^2$] en [$cm^2$]}

\gridspace

\textbf{Exercice 10.} \quad \text{Convertir 27 [$m^3$] en [$cm^3$]}

\gridspace

\textbf{Exercice 11.} \quad \text{Convertir 16 [$kwh$] en [$J$]}

\gridspace

%% ============================================================
%% CHAPITRE 6
%% ============================================================
\section*{Chapitre 6 -- \'{E}nergie}

\rappelbox{green}{%
\begin{itemize}
  \item $E_{p} = mgh$ \quad (\'{e}nergie potentielle de pesanteur)
  \item $E_{c} = \dfrac{1}{2}mv^{2}$ \quad (\'{e}nergie cin\'{e}tique)
  \item $E_{th} = mc\,\Delta T$ \quad (\'{e}nergie thermique, $c = 4180\,\text{J/(kg\ensuremath{\cdot}K)}$ pour l'eau)
  \item $E_{el} = UIt$ \quad (\'{e}nergie \'{e}lectrique)
  \item \textbf{Conservation :} $E_{\text{entr\'{e}e}} = E_{\text{sortie}} + E_{\text{perdue}}$
\end{itemize}
}

\vspace{0.5cm}

\textbf{Exercice 1.} \quad Un objet de $3\,\text{kg}$ est \`{a} $4\,\text{m}$ de hauteur. Calculer $E_{p}$. \; ($g = 9{,}81\,\text{m/s}^{2}$)

\gridspace

\textbf{Exercice 2.} \quad Une bille de $0{,}5\,\text{kg}$ roule \`{a} $6\,\text{m/s}$. Calculer $E_{c}$.

\gridspace

\textbf{Exercice 3.} \quad On chauffe $2\,\text{kg}$ d'eau de $20\,^{\circ}\text{C}$ \`{a} $50\,^{\circ}\text{C}$. Calculer $E_{th}$. \; ($c = 4180\,\text{J/(kg\ensuremath{\cdot}K)}$)

\gridspace

\newpage 

\textbf{Exercice 4.} \quad Une lampe fonctionne sous $12\,\text{V}$ et $0{,}5\,\text{A}$ pendant $5\,\text{min}$. Calculer $E_{el}$.

\gridspace

\textbf{Exercice 5.} \quad Un objet de $2\,\text{kg}$ tombe de $5\,\text{m}$. Quelle est sa vitesse juste avant le sol ? (utiliser $E_{p} = E_{c}$)

\gridspace

\textbf{Exercice 6.} \quad Un ressort ($k = 100\,\text{N/m}$) est comprim\'{e} de $0{,}1\,\text{m}$. Calculer $E_{\text{\'{e}l}} = \dfrac{1}{2}kx^{2}$.

\gridspace

\textbf{Exercice 7.} \quad Une fusée a eau atteint 60m d'altitude. Calculer sa vitesse initial.

\gridspace

\newpage

%% ============================================================
%% CHAPITRE 7
%% ============================================================
\section*{Chapitre 7 -- Puissance et travail}

\rappelbox{orange}{%
\begin{itemize}
  \item $P = \dfrac{E}{t}$ \quad (puissance = \'{e}nergie / temps)
  \item $P = Fv$ \quad (puissance m\'{e}canique)
  \item $W = F \times d \times \cos(\theta)$ \quad (travail d'une force)
  \item $1\,\text{W} = 1\,\text{J/s}$
  \item $1\,\text{kWh} = 3{,}6 \times 10^{6}\,\text{J}$
\end{itemize}
}

\vspace{0.5cm}

\textbf{Exercice 1.} \quad Un moteur fournit $5000\,\text{J}$ en $10\,\text{s}$. Quelle est sa puissance ?

\gridspace

\textbf{Exercice 2.} \quad Une ampoule de $60\,\text{W}$ reste allum\'{e}e $3\,\text{h}$. Quelle \'{e}nergie consomme-t-elle en J ? En kWh ?

\gridspace

\textbf{Exercice 3.} \quad Un ascenseur de $600\,\text{kg}$ monte de $8\,\text{m}$ en $10\,\text{s}$. Quelle puissance ? \; ($g = 9{,}81\,\text{m/s}^{2}$)

\gridspace

\newpage 

\textbf{Exercice 4.} \quad Un cycliste exerce une force de $40\,\text{N}$ et roule \`{a} $7\,\text{m/s}$. Quelle puissance d\'{e}veloppe-t-il ? combien d’énergie pour parcourir 1 km ?

\gridspace

\textbf{Exercice 5.} \quad Convertir $2\,\text{kWh}$ en J.

\gridspace

\textbf{Exercice 6.} \quad Un radiateur de $1500\,\text{W}$ fonctionne $30\,\text{min}$. Quelle \'{e}nergie fournit-il ?

\gridspace


\textbf{Exercice 6.} \quad Une voiture roule à 120 kmh et subit une force de traînée de 400 N, quelle puissance développe t'elle ? combien d'energie pour parcourir 1 km ?

\gridspace

\newpage

%% ============================================================
%% CHAPITRE 8
%% ============================================================
\section*{Chapitre 8 -- Probl\`{e}me int\'{e}gr\'{e} : Mini centrale hydro\'{e}lectrique}

\begin{tcolorbox}[colback=gray!5!white,colframe=black!60,title=\textbf{Situation}]
Un petit barrage retient un r\'{e}servoir d'eau situ\'{e} $h = 50\,\text{m}$ au-dessus d'une turbine. Le r\'{e}servoir contient $V = 2000\,\text{m}^{3}$ d'eau. La turbine a un rendement de $85\,\%$.
\end{tcolorbox}

\vspace{0.5cm}

\textbf{Question 1.} \quad Calculer la masse d'eau dans le r\'{e}servoir. \; ($\rho = 1000\,\text{kg/m}^{3}$)

\gridspace

\textbf{Question 2.} \quad Calculer l'\'{e}nergie potentielle totale de l'eau. \; ($g = 9{,}81\,\text{m/s}^{2}$)

\gridspace

\textbf{Question 3.} \quad En tenant compte du rendement, quelle \'{e}nergie \'{e}lectrique la turbine peut-elle produire ?

\gridspace

\newpage 

\textbf{Question 4.} \quad Convertir cette \'{e}nergie en kWh.

\gridspace

\textbf{Question 5.} \quad Si la turbine a une puissance de $200\,\text{kW}$, combien de temps (en heures) peut-elle fonctionner ?

\gridspace

\textbf{Question 6.} \quad Une locomotive nécessite 2 MW pour accélérer, combien de train peut elle faire partir en meme temps ?

\gridspace


\textbf{Question 7.} \quad Sachant qu'un train transporte 200 Personne et nécessite 5 kwh pour parcourir 1 km, combien de km passager peut déplacer avec l’énergie qui stocké dans le lac ? tiens compte du rendement.

\gridspace

\vspace{1cm}

\begin{center}
\vspace{1cm}
{\Large \textbf{Bravo, tu as termin\'{e} ton cahier de vacances !}}\\[0.3cm]
\end{center}

\end{document}
